\iffalse meta-comment

File: tamu-talk-doc.tex Copyright (C) 2026 Wyatt Smith

It may be distributed and/or modified under the conditions of the
LaTeX Project Public License (LPPL), either version 1.3c of this
license or (at your option) any later version.  The latest version
of this license is in the file

   https://www.latex-project.org/lppl.txt

This file is part of the "tamu-talk bundle" (The Work in LPPL)
and all files in that bundle must be distributed together.

-----------------------------------------------------------------------

The development version of the bundle can be found at

   https://github.com/wsmith-sch/tamu-talk

-----------------------------------------------------------------------

\fi

\DocumentMetadata{pdfversion = 2.0, lang = en, tagging = on, pdfstandard = UA-2}

\documentclass{l3doc}
\usepackage{latex-lab-testphase-l3doc}
\usepackage{unicode-math}

\NewDocumentCommand\opt{m}{\texttt{#1}}

\title{The \pkg{tamu-talk} package: a Texas A\&M theme for \pkg{ltx-talk}}
\author{Wyatt Smith\thanks{E-mail: \href{mailto:216wsmith@gmail.com}{216wsmith@gmail.com}}}
\date{Released 2026-10-24 (v0.1.0)}

\begin{document}

\maketitle

\begin{documentation}

\section{Introduction}

\pkg{tamu-talk} is a lightweight Texas A\&M theme for the \pkg{ltx-talk}
class, a \pkg{beamer}-like presentation class that produces tagged
(accessible) PDF. It restyles the title page, header, and footer.
In addition, it adds the Texas A\&M colors and a few boxes for decorating your presentation.
This package is independently maintained and is not affiliated with,
endorsed by, or sponsored by Texas A\&M University. \pkg{ltx-talk} is
experimental, and the interface of \pkg{tamu-talk} may change with it.
A full example, \texttt{tamu-talk-example.tex}, comes with this
documentation. You are encouraged to copy it as a starting point.

\section{Usage}

\pkg{tamu-talk} works only with the \pkg{ltx-talk} class:
\begin{verbatim}
  \documentclass{ltx-talk}
  \usepackage{tamu-talk}
\end{verbatim}
Options can be given when loading the package, or later in the preamble with
\cs{SetKeys}:
\begin{verbatim}
  \usepackage[header-style = plain]{tamu-talk}
  \SetKeys[tamu-talk]{page-numbers = current}
\end{verbatim}

\section{Obtaining and Using the Texas A\&M Logo}

The Texas A\&M logo is a trademark of Texas A\&M University and is not
included.
Texas A\&M makes its logos available to faculty and staff for marketing and
informational material; anyone else needs the university's approval (see
\url{https://marcomm.tamu.edu/template/logo-downloads/}).

Without the Texas A\&M logo, the title page is laid out without one. If you are
authorized to use it:
\begin{enumerate}
  \item Download the \enquote{TAM Maroon} PNG from the page above, and follow
    the university's logo guidelines.
  \item Save it as \texttt{tamu-logo.png}. Cropping it tightly around the
    logo is encouraged, so that it fills its space on the title page. For the
    sharpest result, export the SVG version to a large PNG and crop that
    instead, for example with
\begin{verbatim}
  rsvg-convert -w 2400 TAM-Logo1.svg \
    | magick - -trim +repage tamu-logo.png
\end{verbatim}
  \item Put it in a \texttt{tamu-talk-logo} folder in your personal texmf
    tree (\texttt{kpsewhich -var-value TEXMFHOME} shows where that is), or
    next to your \texttt{.tex} file. On Overleaf, upload it to your project.
\end{enumerate}

\section{Options}

\begin{description}
  \item[\opt{page-numbers}] Sets the page number in the footer. Values:
    \opt{none} (no number), \opt{current} (the frame number; frames after
    \cs{appendix} are numbered A1, A2, \dots) or \opt{total}
    (\enquote{frame/total}; the total stops at \cs{appendix}, so backup
    frames show numbers past it, such as 11/10). The default value is
    \opt{total}.
  \item[\opt{header-style}] Sets the style of the frame header. Values:
    \opt{banner} (a maroon bar the full height of the header), \opt{ribbon}
    (a thinner maroon bar) or \opt{plain} (maroon text on white).
    The default value is \opt{banner}.
  \item[\opt{title-box}] Puts the title and subtitle in a maroon box on the
    title page; without it, they are printed in maroon. Values: \opt{true}
    or \opt{false}. The default value is \opt{true}.
  \item[\opt{title-bold}] Sets the title, and the frame titles, in bold.
    Values: \opt{true} or \opt{false}. The default value is \opt{false}.
  \item[\opt{title-logo}] Controls the Texas A\&M logo on the title page.
    Values: \opt{auto} (the logo is shown if its file is found), \opt{true}
    (a gray placeholder is drawn, and a warning given, when it is not) or
    \opt{false} (no logo). The default value is \opt{auto}.
  \item[\opt{logo-file}] Sets the logo file, with or without extension.
    Value: a file name. The default value is \opt{tamu-logo}.
\end{description}

\section{The title page}

\cs{maketitle} uses \cs{title}, \cs{subtitle}, \cs{author}, \cs{institute}
and \cs{date}. If \cs{institute} is not given, \enquote{Texas A\&M University,
College Station TX} is used. Its optional argument sets the layout:
\begin{verbatim}
  \maketitle[title-size = 20pt, logo-height = 30mm]
\end{verbatim}
\begin{description}
  \item[\opt{title-size}, \opt{subtitle-size}] Font sizes of the title and
    subtitle. The default values are \opt{18pt} and \opt{11pt}.
  \item[\opt{logo-width}] Width of the logo. The default value is \opt{.25\cs{textwidth}}.
  \item[\opt{logo-height}] Height of the center of the logo above the bottom
    of the slide. The default value is \opt{28mm}.
\end{description}
If a long title would run into the logo, the logo is moved below the text.
You are encouraged to avoid long titles, however.

\section{Colors}

{\raggedright
The Texas A\&M colors are \opt{TAMUMaroon}, \opt{TAMUTan}, \opt{TAMUGray},
\opt{TAMUDarkGray}, \opt{TAMULightGray} and \opt{TAMUNearBlack}. Two off-brand
colors, \opt{forest} and \opt{slate}, pair well with maroon.\par}

\section{Boxes}

Just a few named \pkg{tcolorbox} environments that look good on slides and
use our color palette.
\begin{description}
  \item[\opt{headerbox}\marg{title}] A maroon header bar over a tan body.
  \item[\opt{lasso}\oarg{maroon\textbar{}black}\marg{title}] A border with a
    title tab.
  \item[\opt{sidebar}\oarg{options}] A maroon bar on the left of a tan body,
    without a title. The options are \opt{forest} or \opt{slate} for the
    off-brand colors, or any \pkg{tcolorbox} keys.
  \item[\opt{elevate}\oarg{title}] Rounded corners and a drop shadow, with a
    header bar if a title is given.
  \item[\opt{outline}\oarg{line width}] A thin maroon border only.
\end{description}

\section{Known limits}

\pkg{tamu-talk} uses a few internals of \pkg{ltx-talk} and \pkg{geometry},
so a new release of \pkg{ltx-talk} can break it. Please report problems at
\url{https://github.com/wsmith-sch/tamu-talk/issues}.

\end{documentation}

\end{document}
